\subsubsection{Polygon Area}

\paragraph{Idea intuitiva.}
El area (con signo) de un poligono se calcula con la \textbf{Shoelace formula}: $A = \frac{1}{2} \sum_{i=0}^{n-1} (x_i y_{i+1} - x_{i+1} y_i)$, donde el indice es modulo $n$. El snippet retorna el \textbf{doble} del area (en valor absoluto). La formula sale del ``sumar areas de trapezoides'': cada lado del poligono define un trapezoide con el eje $x$; la suma con signo da el area.

\paragraph{Propiedades clave (invariantes).}
\begin{itemize}\setlength\itemsep{0pt}
	\item El signo indica la orientacion: positivo si counter-clockwise.
	\item Para el doble: no dividir por 2 si no es necesario (evita fracciones).
	\item Funciona para poligonos simples (no auto-intersectantes).
	\item Si el poligono es no convexo, la formula da el area ``con signo''.
\end{itemize}

\paragraph{Codigo clave.}
La suma de trapezoides:
\begin{verbatim}
long long shoelace(const vector<Point>& p) {
    long long s = 0;
    int n = p.size();
    for (int i = 0; i < n; ++i)
        s += (long long)p[i].x * p[(i+1)%n].y - (long long)p[(i+1)%n].x * p[i].y;
    return abs(s);  // doble del area
}
\end{verbatim}

\paragraph{Complejidad.}
\begin{itemize}\setlength\itemsep{0pt}
	\item $O(n)$.
	\item Constante minima.
\end{itemize}

\paragraph{Donde tocar para modificar.}
\begin{itemize}\setlength\itemsep{0pt}
	\item \textbf{Area exacta}: dividir por 2 (con cuidado si da fraccion; usar \texttt{double}).
	\item \textbf{Area firmada}: quitar \texttt{abs} y dejar el signo.
	\item \textbf{Punto dentro de poligono convexo}: con la hull, chequear que esta a un mismo lado de todas las aristas.
	\item \textbf{Triangulacion}: descomponer en triangulos y sumar areas.
\end{itemize}

\paragraph{Trampas y errores frecuentes.}
\begin{itemize}\setlength\itemsep{0pt}
	\item Si los puntos no estan en orden (CCW o CW), el area firmada puede ser incorrecta. Para signed area consistente, usar siempre el mismo orden.
	\item Para poligonos con holes, sumar las areas con signo apropiado.
	\item Con doubles, aniadir tolerancia \texttt{eps}.
\end{itemize}

\paragraph{Codigo.}
Ver \texttt{cpp.tex}, seccion \textit{Polygon Area}.

\vspace{0.5em}
